\emph{फलन}~$f\colon A \to B$ हे \emph{प्रांत}~$A$ च्या प्रत्येक घटकाची \emph{सहप्रांत}~$B$ च्या एका आणि एकाच घटकाशी संगती लावते. $x \in A$ असेल, तर $f$ ने $x$ ची ज्या $y$ शी संगती लावली आहे त्याला \emph{फलसाधक}~$x$ साठी $f$ चे \emph{मूल्य}~$f(x)$ म्हणतो. $A$ हा जोड्यांचा संच असेल, तर $f$ हे दोन फलसाधक घेणारे फलन मानता येते. $f$ ची \emph{व्याप्ती}~$\ran{f}$ म्हणजे $f$ च्या सर्व मूल्यांचा $B$ मधील उपसंच.

$\ran{f} = B$ असेल, तर $f$ ला \emph{आच्छादक} म्हणतात. $f(x)$ हे या अर्थाने एकमेव मूल्य असते की $f$ हे $x$ ची संगती एकाच $f(x)$ शी लावते. कोणत्याही दोन भिन्न फलसाधकांची संगती एकाच मूल्याशी लागत नाही या अर्थानेही $f(x)$ एकमेव असेल, तर $f$ ला \emph{एकास-एक} म्हणतात. जी फलने एकास-एक आणि आच्छादक हे दोन्ही गुणधर्म बाळगतात, त्यांना \emph{एकास-एक व आच्छादक} म्हणतात.

एकास-एक व आच्छादक फलनाचे एकमेव \emph{व्युत्क्रम फलन}~$f^{-1}$ असते. फलने एकापुढे एकही जोडता येतात. $(g \circ f)$ हे $f$ चे $g$ बरोबरचे \emph{संयोजन} आहे. एकास-एक फलनांचे संयोजन एकास-एक असते आणि आच्छादक फलनांचे संयोजन आच्छादक असते. $(f^{-1} \circ f)$ हे एकरूपता फलन आहे.

प्रत्येक $x \in A$ साठी $f$ चे मूल्य असलेच पाहिजे ही अट शिथिल केली तर \emph{अंशतः फलनाची} संकल्पना मिळते. $f\colon A \pto B$ हे अंशतः फलन असेल आणि फलसाधक~$x$ साठी $f$ चे मूल्य असेल, तर $f(x)$ \emph{परिभाषित} आहे असे म्हणतो व $f(x) \fdefined$ असे लिहितो. इतर वेळी $f(x)$ \emph{अपरिभाषित} आहे असे म्हणतो व $f(x) \fundefined$ असे लिहितो. प्रत्येक (अंशतः) फलन~$f$ शी $f$ चा \emph{आलेख}~$R_f$ जोडलेला असतो. $f(x) = y$ असेल तेव्हा आणि केवळ तेव्हाच हा संबंध लागू असतो.
